\section{Growth and development}
Two words are often confused. \kw{Growth} is the increase of the size or the mass of the crop (dry matter, leaf area); \kw{development} is the progress through the stages of the life cycle (germination, emergence, flowering, ripening). Development depends mainly on temperature (and for some crops on the day length), and is measured by the development stage, which Daisy counts from 0 (emergence) to 1 (flowering) to 2 (maturity). Growth depends on the light, the water and the nutrients, and takes place at the rate that the stage allows; the two together decide the yield.

\section{Photosynthesis: C$_3$ and C$_4$ crops}
In \kw{photosynthesis} the leaf takes carbon dioxide from the air through its stomata and uses the energy of light to make sugars. Two pathways are important. In \emph{C$_3$} plants (wheat, rice, barley, potato, sugar beet) the first product has three carbon atoms; at high temperature and light the enzyme also takes oxygen (\emph{photorespiration}), which wastes part of the carbon. In \emph{C$_4$} plants (maize, sorghum, sugar cane) a carbon-concentrating mechanism suppresses the photorespiration, so that they use light, water and nitrogen more efficiently in a warm climate: about 30 per cent more dry matter per unit of water, and a light-saturation that is not reached in full sun. The two groups have different parameters in the light-response function used by the model.

\section{Respiration}
The plant uses a part of the carbon it fixes for its own life. \kw{Maintenance respiration} keeps the existing tissue alive, is proportional to the mass of each organ and rises with the temperature by the $Q_{10}$ law (Equation~\ref{eq:dm-q10} of Chapter~\ref{chap:17}); \kw{growth respiration} is the cost of building new tissue, about 25 to 30 per cent of the carbon used. Roughly half of the carbon fixed by a crop in a season is respired and the rest becomes dry matter, which is why the net production is much smaller than the gross.

\section{Stomata and the exchange of water and carbon dioxide}
The stomata are pores in the leaf; they open to let in carbon dioxide and, as a price, let out water vapour. The opening is controlled by light, by the concentration of carbon dioxide inside the leaf, by the humidity of the air and by the water status of the plant. Under stress the stomata close, the transpiration falls and the photosynthesis falls too; the ratio of carbon gained to water lost is the \kw{water-use efficiency} (WUE). The Ball--Berry model used in the crop chain of this chapter links the stomatal conductance to the photosynthesis, the humidity and the carbon dioxide at the leaf surface.

\section{Partitioning and the structure of the crop}
The carbon fixed is divided among the organs (roots, stems, leaves, storage organs) by \kw{partitioning} fractions that change with the development stage: early, most goes to the roots and leaves, so that the crop can capture light and water; after flowering, most goes to the storage organs (the grain, the tuber). The leaf area index, the root depth and the height are consequences of the dry matter of the organs, through specific leaf area, root-length density and the allometry of the stem. The harvest index (the share of the dry matter in the grain) is the final result of the partitioning.

\section{Management of a field and its model}
\kw{Management} is the set of operations that a farmer applies: sowing, tillage, fertilising, irrigating, spraying, harvesting. In a model each operation is an \emph{action} that is applied at a date or when a \emph{condition} is met (a soil water suction beyond a trigger, a temperature sum, a development stage). A list of actions in order, with their conditions, is a \emph{management program}. The trigger rule of Chapter~\ref{chap:03} is an example of a conditional action. The management program is the part of the model that the farmer controls and that the study varies, so that its description must be complete: an action not in the program does not happen.
